% ============================================================================
% Fichier  : partie0/P00_tiroir.tex
% Titre    : Le schéma directeur dort dans un tiroir
% Auteur   : MINKA MI NGUIDJOÏ Thierry Emmanuel
% Version  : 1.0 -- novembre 2026
% Corrigés : corriges/Corr_P00.tex
% ============================================================================

\chapter{Le schéma directeur dort dans un tiroir}
\label{chap:prologue}
\chaptermark{Prologue}

\epigraph{%
    « Un plan que personne ne suit n'est pas un mauvais plan.\\
    C'est un document. »%
}{Principe du cours}

\niveauChapitre{Introductif}{%
    Aucun prérequis
}

\begin{boxobjectifs}
\begin{itemize}
    \item Reconnaître un schéma directeur qui ne sert pas.
    \item Poser les trois questions du cours~: aligné, finançable, gouverné.
    \item Coter une proposition \OUI{}, \NON{} ou \NEK{} avec une pièce.
    \item Situer les trois organisations du cours.
\end{itemize}
\end{boxobjectifs}

\bigskip

% ============================================================================
\section*{Avant de commencer : « nous en avons déjà un »}
% ============================================================================

Yaoundé, mardi 14~janvier~2025. L'Agence nationale de paiement des prestations\footnote{L'ANPP, ses dirigeants,
ses projets et ses chiffres sont fictifs, comme LiZbiZ, la BANQUE SANAGA et le cabinet AUDIX CONSEIL. Ces
organisations figurent dans d'autres cours de l'auteur.}, l'ANPP, verse chaque mois des prestations à
186\,000~bénéficiaires par dix antennes régionales. Elle emploie 412~agents. Son conseil d'administration vient de
commander un schéma directeur pour 2025-2029 au cabinet AUDIX CONSEIL.

À la réunion de lancement, la cheffe de la division informatique, Mme~NGO~BIYONG, pose un classeur sur la
table. « Nous en avons déjà un. Il couvrait 2019 à 2023. Je suis arrivée en 2021 et personne ne m'en a parlé. »
Le directeur général, M.~OWONA, ajoute qu'en mars~2024 une panne a arrêté les paiements pendant quatre jours.
Le projet e-RECETTES, mis en service en 2023, n'avait pas de site de secours. « Le conseil veut un nouveau
document. Je voudrais surtout qu'il serve. »

M.~BEYALA, chef de mission, confie le classeur à Mme~EKANI et au consultant junior. Leur première tâche n'est
pas d'écrire. Elle est de comprendre pourquoi le document précédent n'a servi à rien.

\begin{boxalert}
\textbf{Violation typique.} Commander un nouveau schéma directeur sans avoir fait le bilan du précédent.

\smallskip
\textbf{Règle~:} les causes qui ont laissé dormir un document le laisseront dormir une seconde fois. Le bilan
du schéma directeur échu est la première pièce du nouveau.
\end{boxalert}

% ============================================================================
\section{Ce que contenait le schéma directeur 2019-2023}
\label{sec:p-contenu}
% ============================================================================

Le document compte 214~pages. Il a été rédigé en 2018 par un cabinet, validé par le conseil d'administration le
17~décembre~2018 et payé 86~millions de francs CFA. Mme~EKANI en dresse l'inventaire en une matinée.

\begin{tableaucours}{L{5.2cm}cL{5.6cm}}{%
    \tableauHeader{Rubrique} &
    \tableauHeader{Pages} &
    \tableauHeader{Ce que l'on y trouve}
}
État des lieux & 118 & Inventaire du matériel, poste par poste \\
Orientations & 9 & Sept axes, sans lien avec le plan stratégique de l'Agence \\
Projets & 61 & 23 fiches de projet, sans priorité ni responsable \\
Budget & 4 & Un total de 6,3~milliards, sans échéancier \\
Gouvernance & 2 & Un comité de pilotage, jamais réuni \\
Annexes & 20 & Glossaire, listes \\
\end{tableaucours}
\vspace{-0.8em}
\begin{center}{\small\itshape Tableau~P.1. Anatomie du schéma directeur 2019-2023 de l'ANPP}\end{center}

Plus de la moitié du document décrit du matériel qui a été remplacé depuis. Les orientations tiennent en neuf
pages. Le budget tient en quatre, pour un montant qui représente dix fois le budget informatique annuel de
l'Agence. Aucune ligne ne dit d'où viendrait l'argent.

% ============================================================================
\section{Ce qui s'est passé en cinq ans}
\label{sec:p-bilan}
% ============================================================================

\begin{boxscenario}{Mission ANPP. Le bilan du schéma directeur 2019-2023, 21 janvier 2025}
\textbf{Constatations.}
\begin{itemize}
    \item Sur les 23 projets inscrits, 6 ont été livrés, 4 ont été commencés puis arrêtés, 13 n'ont jamais
          démarré.
    \item Sur la même période, l'Agence a mené 11 projets qui ne figuraient pas au schéma directeur, dont
          e-RECETTES, le plus coûteux.
    \item Le comité de pilotage prévu ne s'est jamais réuni. Aucun procès-verbal n'existe.
    \item Le conseil d'administration n'a reçu aucun point d'avancement après la validation.
    \item Les dépenses informatiques réelles s'élèvent à 3,1~milliards sur cinq ans, soit la moitié du montant
          annoncé.
\end{itemize}

\textbf{Analyse} (confiance élevée sur les chiffres, moyenne sur les causes). Le document n'a pas guidé les
décisions. Les projets réalisés ont répondu à des urgences ou à des financements extérieurs. Personne n'était
chargé de comparer ce qui se faisait à ce qui était prévu.
\end{boxscenario}

M.~BEYALA écrit trois questions au tableau. Elles servent de fil à tout le cours.

\begin{enumerate}
    \item \textbf{Aligné~?} À quelle stratégie de l'Agence ce document répondait-il~?
    \item \textbf{Finançable~?} Avec quel argent, inscrit dans quel budget, engagé selon quelle procédure~?
    \item \textbf{Gouverné~?} Qui décidait, qui suivait, qui rendait compte~?
\end{enumerate}

% ============================================================================
\section{Huit premières propositions}
\label{sec:p-lecture}
% ============================================================================

Avant tout diagnostic, l'équipe cote huit propositions simples sur le document. Chacune reçoit une cote et une
pièce. C'est la convention de tout le cours.

\begin{boxlecture}{Le schéma directeur 2019-2023, 21 janvier 2025}
\begin{tableaucours}{L{7.6cm}cL{3.4cm}}{%
    \tableauHeader{Proposition} &
    \tableauHeader{Cote} &
    \tableauHeader{Preuve}
}
Le document cite le plan stratégique de l'Agence & \NON{} & Aucune mention \\
Chaque projet est rattaché à un objectif de l'Agence & \NON{} & 23 fiches lues \\
Les projets sont classés par priorité & \NON{} & Liste alphabétique \\
Le budget est réparti par année & \NON{} & Total unique \\
Une source de financement est indiquée & \NON{} & Aucune \\
Une instance de pilotage est désignée & \OUI{} & Page 198 \\
Cette instance s'est réunie & \NON{} & Aucun procès-verbal \\
Le conseil a suivi l'exécution & \NEK{} & Procès-verbaux de 2020 non retrouvés \\
\end{tableaucours}

\medskip
\textbf{Décompte.} Un \OUI{}, six \NON{}, un \NEK{} sur huit.
\hfill\lecture{0{,}125}{0{,}125}{1{,}06}

\medskip
\textbf{Lecture.} Le document n'est ni aligné ni finançable. Il désigne une instance, ce qui vaut un \OUI{}, mais
l'instance n'a pas fonctionné. Le \NEK{} est de preuve~: les procès-verbaux du conseil pour 2020 manquent aux
archives. Avec $U = 0{,}125$, l'incertitude ne gêne pas la conclusion. Six écarts sont établis.
\end{boxlecture}

\begin{boxdef}
\textbf{Schéma directeur.} Document qui fixe, pour une durée de trois à cinq ans, la cible du système
d'information d'une organisation, les projets qui y conduisent, leur ordre, leur coût et la manière dont
l'ensemble sera piloté. Il découle de la stratégie de l'organisation. Il n'est ni un inventaire ni un catalogue
de souhaits.
\end{boxdef}

% ============================================================================
\section{Le plan du cours}
\label{sec:p-plan}
% ============================================================================

Le cours suit l'ordre de la mission. Les fondements donnent le vocabulaire et les théories
(Chapitres~\ref{chap:vocabulaire} à~\ref{chap:alignement}). Le contexte international puis national dit ce qui
s'impose à l'Agence avant qu'elle ne décide (Chapitres~\ref{chap:normes} à~\ref{chap:droit}). La méthode
déroule les cinq temps d'un schéma directeur (Chapitres~\ref{chap:cadrer} à~\ref{chap:financer}). La mise en
œuvre traite le pilotage et la maîtrise des risques (Chapitres~\ref{chap:piloter} et~\ref{chap:maitriser}).
Les deux derniers chapitres ouvrent sur les données, l'intelligence artificielle et l'audit du schéma directeur
lui-même.

\begin{boxmetier}{Consultant}
Lire le document précédent avant toute réunion. Compter ce qui a été fait, ce qui ne l'a pas été et ce qui s'est fait en dehors.
\end{boxmetier}

\begin{boxmetier}{Directeur des systèmes d'information}
Garder la trace des décisions. Un projet hors schéma directeur n'est pas une faute. Un projet dont personne ne sait pourquoi il a été lancé en est une.
\end{boxmetier}

\begin{boxmetier}{Dirigeant}
Demander chaque année où en est le schéma directeur. Un document que le dirigeant ne réclame pas n'est suivi par personne.
\end{boxmetier}

\begin{boxresume}
\begin{itemize}
    \item Un schéma directeur fixe une cible, des projets ordonnés, un coût et un pilotage. Il découle de la stratégie.
    \item Trois questions le jugent~: est-il aligné, finançable, gouverné~?
    \item À l'ANPP, 6 projets sur 23 ont été livrés et 11 projets se sont faits hors du document.
    \item Chaque proposition se cote \OUI{}, \NON{} ou \NEK{}, avec une pièce.
    \item Le bilan du document échu ouvre le nouveau.
\end{itemize}
\end{boxresume}

\begin{boxinfo}
\textbf{Points de contrôle.} Vous avez compris ce chapitre si vous savez~:
\begin{itemize}
    \item définir un schéma directeur et dire ce qu'il n'est pas~;
    \item énoncer les trois questions du cours~;
    \item expliquer la différence entre un \NON{} et un \NEK{}.
\end{itemize}
\end{boxinfo}

\begin{boxexo}[title={\ding{45}\quad Exercice~P.1. Le bilan en chiffres \niveaubase}]
Calculez, pour le schéma directeur 2019-2023 de l'ANPP, le taux de projets livrés, le taux de projets hors document parmi les projets menés et le taux d'exécution du budget.
\end{boxexo}

\begin{boxexo}[title={\ding{45}\quad Exercice~P.2. Coter votre organisation \niveaumoyen}]
Prenez le schéma directeur d'une organisation que vous connaissez, ou à défaut son plan informatique. Cotez les huit propositions du prologue. Calculez $\mu$ et $U$. Rédigez la lecture en cinq lignes.
\end{boxexo}

\begin{boxexo}[title={\ding{45}\quad Exercice~P.3. La lettre au conseil \niveauexpert}]
Rédigez la note d'une page que M.~BEYALA adresse au président du conseil, M.~TCHUENTE, pour expliquer pourquoi la mission commence par un bilan. Séparez constatations et analyse.
\end{boxexo}

\bigskip

\begin{boxinfo}
\textbf{Pour aller plus loin.}
\begin{itemize}
    \item Le chapitre sur le schéma directeur dans Peppard et Ward~\cite{wardpeppard2016}.
    \item L'ISO/IEC~38500~\cite{iso38500}, pour le rôle de l'organe de gouvernance.
\end{itemize}
\end{boxinfo}

\begin{boxinfo}
\textbf{Transition.} Avant de reconstruire, il faut des mots exacts. Stratégie, système d'information, numérique, plan, schéma directeur, feuille de route~: le Chapitre~\ref{chap:vocabulaire} les sépare.
\end{boxinfo}
